% 摘要
\begin{abstract}
\noindent 这篇文章收拢了从 1950 年代资产定价计量到 2026 年大模型交易智能体之间，关于“人工智能如何被用进金融交易”的论文、预印本、官方仓库与基准报告。我不打算把它写成一份方法清单，而是先立一个能把整条链路串起来的命题：交易的本质是“信息 $\to$ 定价偏差 $\to$ 可交易缝隙”的识别与收割，AI 的每一种应用，无非是在这条链条的某一环替换或增强了人的表示学习、预测或决策能力；因此它的能力上限，由“信息存量的天花板”和“执行摩擦的地板”共同夹住。这个命题解释了一件长期被回避的事——为什么论文里动辄年化几十个百分点的策略，一旦落到真实账户上就大面积缩水。

正文用一张“方法族 $\times$ 交易环节”的双轴框架来组织。纵轴是六个方法族：监督学习与 alpha 挖掘、深度学习与序列建模、强化学习、自然语言处理与大模型、生成式模型、图神经网络；横轴是五个交易环节：信号生成、组合构建、交易执行、风险管理、另类数据与微观结构。每一族都按“问题设定 $\to$ 原理 $\to$ 落地要素 $\to$ 复现状态 $\to$ 失效边界”走一遍，而不是停在“某方法在某数据集上取得了更好指标”。

落地是我写得最用力的部分。第 \ref{sec:evidence} 章的\textbf{证据表}收录了超过 80 条核心来源，全部链回 DOI / arXiv / 官方仓库；第 \ref{sec:matrix} 章的\textbf{统一比较矩阵}把“数据、样本区间、任务、方法、样本外检验、交易成本口径、评测协议、开源许可”逐项对齐，好让跨论文的数字有可比的基础；第 \ref{sec:audit} 章专门审计评测污染——幸存者偏差、前瞻偏差、数据窥探、多重检验、回测过拟合；第 \ref{sec:repro_audit} 章的\textbf{复现状态核验表}把 Qlib、FinRL 等公开基准和一批代表性工作逐一标注为“可复现 / 部分复现 / 不可复现”并给出原因；第 \ref{sec:repro_protocol} 章给出一套按 P0--P3 分级的可复现实验方案，配套的参考实现骨架放在随附的 \texttt{code/} 目录里。结论附置信度分级，边界条件如实写明。

\end{abstract}
